\documentclass[11pt,a4paper]{article}

\usepackage[utf8]{inputenc}
\usepackage[T1]{fontenc}
\usepackage[italian]{babel}
\usepackage{amsmath,amssymb,amsthm}
\usepackage{bm}
\usepackage[margin=2.4cm]{geometry}
\usepackage{enumitem}
\usepackage{xcolor}
\usepackage{mdframed}
\usepackage{hyperref}
\hypersetup{colorlinks=true,linkcolor=blue!55!black,urlcolor=blue!55!black}

\newtheorem{theorem}{Teorema}
\newtheorem{lemma}{Lemma}
\newtheorem{proposition}{Proposizione}
\theoremstyle{definition}
\newtheorem{definition}{Definizione}
\newtheorem{remark}{Osservazione}

\newmdenv[linecolor=blue!40!black,linewidth=0.6pt,
  backgroundcolor=blue!4,roundcorner=4pt,
  innertopmargin=6pt,innerbottommargin=6pt]{keybox}

\newcommand{\w}{\bm{w}}
\newcommand{\X}{\bm{X}}
\newcommand{\y}{\bm{y}}
\newcommand{\A}{\bm{A}}
\newcommand{\bvec}{\bm{b}}
\newcommand{\Q}{\bm{Q}}
\newcommand{\Hess}{\bm{H}}
\newcommand{\Id}{\bm{I}}
\newcommand{\wstar}{\w^{*}}
\newcommand{\Mmap}{\mathcal{M}}
\newcommand{\diag}{\operatorname{diag}}
\newcommand{\sign}{\operatorname{sign}}
\newcommand{\spec}{\rho}
\newcommand{\lam}{\lambda}
\newcommand{\eps}{\varepsilon}

\title{\textbf{IRLS: rate di convergenza, spiegato passo passo}\\[2pt]
\large La Jacobiana $\nabla\Mmap(\wstar)=\Id-\Q(\wstar)^{-1}\nabla^2 f_\eps(\wstar)$
e la dimostrazione di $\spec(\nabla\Mmap(\wstar))<1$}
\author{Note di studio --- progetto CM (ELM + LASSO, algoritmo A1)}
\date{\today}

\begin{document}
\maketitle

\begin{abstract}
\noindent Documento di supporto allo studio (per me e per Matthew). Ricostruisce, partendo
dalle basi, perché IRLS converge \emph{linearmente} sul problema ELM\,+\,LASSO e perché
il rate è \emph{lento} (fattore di contrazione vicino a $1$). Il cuore sono due cose:
(i)~la forma della Jacobiana della mappa di IRLS nel punto fisso, e (ii)~la dimostrazione
che il suo raggio spettrale è strettamente minore di $1$. Tutto è coerente con il
paragrafo \emph{Convergence rate} del capitolo IRLS del report; qui è solo svolto con più
dettaglio e senza passaggi impliciti.
\end{abstract}

\tableofcontents

% ===================================================================
\section{Il problema e cosa fa IRLS}

Il problema che vogliamo risolvere (ELM con regolarizzazione LASSO) è
\begin{equation}
\label{eq:problem}
\min_{\w\in\mathbb{R}^H}\; f(\w)=\tfrac12\|\X\w-\y\|_2^2+\lam\|\w\|_1 ,
\end{equation}
con $\X\in\mathbb{R}^{M\times H}$ la matrice delle feature dello strato nascosto (assunta a
\emph{rango colonna pieno}, $M\ge H$) e $\lam>0$ il peso LASSO. Il termine $\|\w\|_1$ rende
$f$ \textbf{non differenziabile} dove qualche componente è zero, quindi non si possono usare
metodi a gradiente direttamente.

\medskip
\noindent\textbf{Idea di IRLS.} A ogni passo si sostituisce $|w_i|$ con un maggiorante
quadratico (liscio) costruito attorno al punto corrente $\w_k$, e si minimizza il modello
risultante. Concretamente, usando la disuguaglianza $(|w_i|-|(w_k)_i|)^2\ge 0$ si ottiene
\[
|w_i|\;\le\;\frac{w_i^2}{2|(w_k)_i|}+\frac{|(w_k)_i|}{2}\qquad((w_k)_i\neq 0),
\]
e sommando su $i$ il termine $L_1$ viene maggiorato da una forma quadratica pesata. Il modello
(surrogato) da minimizzare a ogni iterazione è quindi un \textbf{minimi quadrati pesati con
regolarizzazione $L_2$}, che ha soluzione in forma chiusa risolvendo un sistema lineare SPD.

\medskip
\noindent\textbf{Pesi con soglia.} Per evitare la divisione per zero quando una componente si
avvicina a $0$, i pesi usano una soglia $\eps>0$:
\begin{equation}
\label{eq:weights}
(\bm{W}_k)_{ii}=\Big(\max\!\big(|(w_k)_i|,\eps\big)\Big)^{-1/2}.
\end{equation}
Questa soglia, come vedremo, è ciò che trasforma il problema risolto da IRLS in una versione
\emph{smussata} del LASSO, ed è anche il motivo per cui la teoria diventa pulita.

% ===================================================================
\section{Cosa significa ``rate'': lineare vs sublineare}

Chiamiamo $g_k$ una misura dell'errore al passo $k$ (per esempio $\|\w_k-\wstar\|$, oppure
il gap $f(\w_k)-f^\star$).

\begin{definition}[Convergenza lineare]
La successione converge \emph{linearmente} con fattore $\spec\in(0,1)$ se
\[
\lim_{k\to\infty}\frac{g_{k+1}}{g_k}=\spec,\qquad 0<\spec<1 .
\]
A ogni passo l'errore viene moltiplicato (asintoticamente) per la stessa costante $\spec$.
\end{definition}

\begin{definition}[Convergenza sublineare]
La convergenza è \emph{sublineare} se $g_{k+1}/g_k\to 1$ (l'errore cala, ma il guadagno per
passo svanisce). Esempio tipico: il metodo del subgradiente, con $g_k\sim 1/\sqrt{k}$.
\end{definition}

\begin{keybox}
\textbf{Lettura su grafico semilog} (asse $y=\log g_k$, asse $x=k$):
\begin{itemize}[nosep]
\item lineare $\Rightarrow$ \emph{retta} di pendenza $\log\spec$. Se $\spec\approx0.99$ la
  retta è quasi orizzontale (lineare ma \emph{lenta});
\item sublineare $\Rightarrow$ curva \emph{concava} che si appiattisce progressivamente.
\end{itemize}
Una lineare con $\spec\approx0.99$ su poche iterazioni è facilmente scambiabile a occhio
per una sublineare: è esattamente il malinteso da chiarire nella figura del report.
\end{keybox}

% ===================================================================
\section{IRLS come iterazione a punto fisso}
\label{sec:fixedpoint}

La soluzione del sotto-problema di IRLS (equazione normale pesata) è, a meno della costante,
\begin{equation}
\label{eq:system}
\Q_k\,\w_{k+1}=\bvec,\qquad
\Q_k=\A+\lam\,\diag\!\Big(\tfrac{1}{\max(|(w_k)_i|,\eps)}\Big),\quad
\A=\X^{\!\top}\X,\ \ \bvec=\X^{\!\top}\y .
\end{equation}
(La costante $\lam$ qui è il peso LASSO; nel report compare anche $\lam_{\text{IRLS}}=\tfrac12\lam_{\text{LASSO}}$
dovuto al fattore $1/2$ del maggiorante, ma nella matrice finale i conti tornano a $\lam_{\text{LASSO}}$.)

Poiché $\Q_k$ dipende \emph{solo} da $\w_k$, l'intera iterazione è una funzione del punto corrente:
\begin{equation}
\label{eq:map}
\boxed{\ \w_{k+1}=\Mmap(\w_k),\qquad \Mmap(\w)=\Q(\w)^{-1}\bvec\ }
\end{equation}
con $\Q(\w)=\A+\lam\diag\!\big(1/\max(|w_i|,\eps)\big)$. La soluzione $\wstar$ è un
\textbf{punto fisso}: $\Mmap(\wstar)=\wstar$, cioè $\Q(\wstar)\wstar=\bvec$.

\begin{keybox}
\textbf{Principio guida.} Studiare la velocità di IRLS = studiare quanto la mappa $\Mmap$
\emph{contrae} le distanze vicino a $\wstar$. Se vicino al punto fisso $\Mmap$ avvicina i
punti a $\wstar$ di un fattore $\spec<1$ a ogni applicazione, allora la convergenza è
lineare con rate $\spec$. Questo fattore è il raggio spettrale della Jacobiana di $\Mmap$
in $\wstar$.
\end{keybox}

% ===================================================================
\section{Il problema ``vero'' che IRLS risolve: la versione smussata}
\label{sec:smoothed}

La soglia $\eps$ in \eqref{eq:weights} non è solo un trucco numerico: cambia la funzione
obiettivo. Cerchiamo la funzione liscia $\psi_\eps$ il cui maggiorante quadratico ha
esattamente i pesi $1/\max(|w_i|,\eps)$. Imponendo che la derivata di $\psi_\eps$ sia
$\lam\,w_i/\max(|w_i|,\eps)$ si ottiene
\begin{equation}
\label{eq:huber}
\psi_\eps(t)=
\begin{cases}
|t|, & |t|>\eps,\\[2pt]
\dfrac{t^2}{2\eps}+\dfrac{\eps}{2}, & |t|\le\eps,
\end{cases}
\end{equation}
cioè una smussatura di tipo Huber del valore assoluto (lineare lontano da $0$, quadratica
in un intorno di $0$, raccordata in modo $C^1$). Il problema effettivamente risolto da IRLS è
\begin{equation}
\label{eq:feps}
\min_{\w}\ f_\eps(\w)=\tfrac12\|\X\w-\y\|_2^2+\lam\sum_{i=1}^H\psi_\eps(w_i).
\end{equation}
Il minimo di $f_\eps$ dista dal minimo del LASSO originale di $O(\eps)$ in valore obiettivo,
e lo recupera per $\eps\to0$. Con $\eps=10^{-8}$ la differenza è sotto la soglia di
rilevabilità rispetto a un solver di riferimento.

\begin{proposition}[Forte convessità]
\label{prop:strong}
Se $\X$ ha rango colonna pieno, allora $\A=\X^{\!\top}\X\succ0$ e $f_\eps$ è
\emph{fortemente convessa}, quindi ha un unico minimizzatore $\wstar$.
\end{proposition}
\begin{proof}
Per ogni $\bm v\neq\bm0$ si ha $\bm v^{\!\top}\A\bm v=\|\X\bm v\|_2^2>0$ (rango pieno),
dunque $\A\succ0$. La parte di penalità $\sum_i\psi_\eps(w_i)$ è convessa ($\psi_\eps$ è
convessa: derivata seconda $0$ o $1/\eps$, mai negativa). Somma di una funzione fortemente
convessa e di una convessa è fortemente convessa; l'unicità del minimo segue.
\end{proof}

\noindent La Hessiana del problema smussato è
\begin{equation}
\label{eq:hessian}
\nabla^2 f_\eps(\w)=\A+\lam\,\diag\!\big(\psi_\eps''(w_i)\big),\qquad
\psi_\eps''(t)=
\begin{cases}
0, & |t|>\eps \quad(\text{tratto lineare}),\\[2pt]
1/\eps, & |t|<\eps \quad(\text{tratto quadratico}).
\end{cases}
\end{equation}
È qui che si gioca tutto: \textbf{sulle componenti attive} ($|w_i|>\eps$) la funzione vera
$|w_i|$ è una retta, curvatura nulla; \textbf{sulle componenti schiacciate a zero}
($|w_i|<\eps$, ``pinned'') la curvatura è $1/\eps$.

% ===================================================================
\section{La Jacobiana della mappa: derivazione da zero}
\label{sec:jacobian}

Vogliamo $\nabla\Mmap(\wstar)$. Partiamo da $\Mmap(\w)=\Q(\w)^{-1}\bvec$ e differenziamo.
Usiamo l'identità $d(\bm A^{-1})=-\bm A^{-1}(d\bm A)\bm A^{-1}$:
\begin{equation}
\label{eq:dM1}
d\Mmap=-\Q^{-1}(d\Q)\,\Q^{-1}\bvec=-\Q^{-1}(d\Q)\,\w ,
\end{equation}
dove nell'ultimo passaggio abbiamo usato che nel punto fisso $\Q(\wstar)^{-1}\bvec=\wstar$.

\medskip
\noindent\textbf{Calcolo di $d\Q$.} La matrice $\Q(\w)$ è $\A$ (costante) più la diagonale
$\lam\,\diag(1/\max(|w_i|,\eps))$. La derivata dell'$i$-esimo elemento diagonale rispetto a
$w_i$:
\[
\frac{d}{dw_i}\,\frac{1}{\max(|w_i|,\eps)}=
\begin{cases}
-\dfrac{\sign(w_i)}{w_i^2}, & |w_i|>\eps,\\[4pt]
0, & |w_i|<\eps.
\end{cases}
\]
Quindi $d\Q$ è diagonale, con entrate non nulle solo sull'insieme attivo.

\medskip
\noindent\textbf{Il prodotto $(d\Q)\,\wstar$.} Sull'insieme attivo l'$i$-esima componente è
\[
\big[(d\Q)\,\wstar\big]_i=\lam\left(-\frac{\sign(w^*_i)}{(w^*_i)^2}\right)(w^*_i)\,dw_i
=-\,\lam\,\frac{\sign(w^*_i)}{w^*_i}\,dw_i
=-\,\frac{\lam}{|w^*_i|}\,dw_i ,
\]
mentre sull'insieme pinned è $0$. In forma matriciale
\begin{equation}
\label{eq:dQw}
(d\Q)\,\wstar=-\underbrace{\lam\,\diag\!\Big(\tfrac{1}{|w^*_i|}\ \text{(attive)},\ 0\ \text{(pinned)}\Big)}_{=\,\Q(\wstar)-\nabla^2 f_\eps(\wstar)}\;d\w .
\end{equation}
L'uguaglianza con $\Q(\wstar)-\nabla^2 f_\eps(\wstar)$ si verifica confrontando \eqref{eq:system}
con \eqref{eq:hessian}:
\[
\Q(\wstar)-\nabla^2 f_\eps(\wstar)=\lam\,\diag\!\Big(\tfrac{1}{\max(|w^*_i|,\eps)}-\psi_\eps''(w^*_i)\Big)
=\lam\,\diag\!\Big(\underbrace{\tfrac{1}{|w^*_i|}-0}_{\text{attive}},\ \underbrace{\tfrac1\eps-\tfrac1\eps}_{\text{pinned}}\Big).
\]
Sostituendo \eqref{eq:dQw} in \eqref{eq:dM1}:
\[
d\Mmap=-\Q^{-1}\Big(-(\Q-\nabla^2 f_\eps)\,d\w\Big)
=\Q^{-1}(\Q-\nabla^2 f_\eps)\,d\w=(\Id-\Q^{-1}\nabla^2 f_\eps)\,d\w .
\]

\begin{keybox}
\begin{equation}
\label{eq:jac}
\boxed{\ \nabla\Mmap(\wstar)=\Id-\Q(\wstar)^{-1}\,\nabla^2 f_\eps(\wstar)\ }
\end{equation}
È \emph{esattamente} la matrice di contrazione di una iterazione Majorization--Minimization.
Non l'abbiamo presa per citazione: l'abbiamo \emph{derivata} differenziando la mappa.
\end{keybox}

\begin{remark}[Quando $\Mmap$ è derivabile]
Il calcolo sopra usa $\psi_\eps''$, che ha un salto in $|w_i|=\eps$. La derivazione è quindi
lecita finché \emph{nessuna} componente di $\wstar$ giace esattamente sul bordo
$|w^*_i|=\eps$ (le attive stanno strettamente sopra, le pinned strettamente sotto). In quel
caso $\Mmap$ e $f_\eps$ sono $C^1$ in $\wstar$. È un'ipotesi tecnica generica (insieme di
misura nulla), ma è un'\emph{assunzione}, non un teorema --- vedi la
sezione~\ref{sec:caveats}.
\end{remark}

% ===================================================================
\section{Dimostrazione che il raggio spettrale e' minore di 1}
\label{sec:proof}

Questo è il risultato centrale. Poniamo per brevità
\[
\bm B:=\Q(\wstar),\qquad \Hess:=\nabla^2 f_\eps(\wstar),\qquad
\nabla\Mmap(\wstar)=\Id-\bm B^{-1}\Hess .
\]

\paragraph{Passo 0 --- entrambe SPD.}
$\bm B=\A+\lam\,\diag(1/\max(|w^*_i|,\eps))$: somma di $\A\succ0$ e di una diagonale
strettamente positiva, quindi $\bm B\succ0$.
$\Hess=\A+\lam\,\diag(\psi_\eps''\ge0)$: somma di $\A\succ0$ e di una diagonale $\succeq0$,
quindi $\Hess\succ0$. Entrambe sono simmetriche definite positive.

\paragraph{Passo 1 --- gli autovalori di $\bm B^{-1}\Hess$ sono reali e positivi.}
$\bm B^{-1}\Hess$ in generale non è simmetrica, ma è \emph{simile} a una matrice simmetrica:
\[
\bm B^{-1}\Hess=\bm B^{-1/2}\big(\bm B^{-1/2}\Hess\,\bm B^{-1/2}\big)\bm B^{1/2},
\qquad
\bm S:=\bm B^{-1/2}\Hess\,\bm B^{-1/2}.
\]
$\bm S$ è simmetrica e, essendo $\Hess\succ0$, è anche definita positiva. Matrici simili hanno
gli stessi autovalori, dunque gli autovalori di $\bm B^{-1}\Hess$ coincidono con quelli di
$\bm S$: sono \textbf{reali e $>0$}. Chiamiamoli $\mu_1\ge\cdots\ge\mu_H>0$.

\paragraph{Passo 2 --- limite superiore: $\mu_i\le 1$.}
Vale $\bm S\preceq\Id$ se e solo se $\Hess\preceq\bm B$, cioè $\bm B-\Hess\succeq0$.
Ma dalla sezione~\ref{sec:jacobian},
\begin{equation}
\label{eq:diff}
\bm B-\Hess=\lam\,\diag\!\Big(\underbrace{\tfrac{1}{|w^*_i|}}_{\text{attive}>0},\ \underbrace{0}_{\text{pinned}}\Big)\ \succeq\ 0 ,
\end{equation}
una matrice diagonale con entrate $\ge0$, quindi semidefinita positiva. Dunque $\Hess\preceq\bm B$,
quindi $\bm S\preceq\Id$, quindi $\mu_i\le1$ per ogni $i$.

\paragraph{Passo 3 --- conclusione.}
Combinando i passi 1--2: $\mu_i\in(0,1]$. Gli autovalori di
$\nabla\Mmap(\wstar)=\Id-\bm B^{-1}\Hess$ sono allora $1-\mu_i\in[0,1)$. Il raggio spettrale è
il massimo modulo:
\[
\spec(\nabla\Mmap(\wstar))=\max_i|1-\mu_i|=1-\mu_{\min}<1,
\]
dove $\mu_{\min}=\mu_H>0$ è il più piccolo autovalore di $\bm B^{-1}\Hess$. Poiché
$\mu_{\min}>0$ \emph{strettamente}, si ha $\spec<1$ strettamente. \qed

\begin{keybox}
\textbf{Conseguenza (teorema di Ostrowski).} Per un'iterazione a punto fisso
$\w_{k+1}=\Mmap(\w_k)$ con $\Mmap$ di classe $C^1$ in $\wstar$ e
$\spec(\nabla\Mmap(\wstar))<1$, il punto $\wstar$ è un \emph{punto di attrazione}: esiste un
intorno da cui IRLS converge a $\wstar$, e la convergenza è \textbf{lineare} con rate
asintotico esattamente $\spec(\nabla\Mmap(\wstar))$. (Ostrowski, in
Ortega--Rheinboldt, \emph{Iterative Solution of Nonlinear Equations}.)
\end{keybox}

% ===================================================================
\section{Perché il rate è vicino a 1 (cioè: perché IRLS è lento)}
\label{sec:slow}

Abbiamo $\spec=1-\mu_{\min}$, quindi ``vicino a $1$'' equivale a ``$\mu_{\min}$ piccolo''.
Il più piccolo autovalore generalizzato è
\[
\mu_{\min}=\min_{\bm v\neq\bm0}\frac{\bm v^{\!\top}\Hess\,\bm v}{\bm v^{\!\top}\bm B\,\bm v}.
\]
Consideriamo una direzione $\bm v$ concentrata su una componente attiva con $|w^*_i|$ piccolo.
Su quella componente $\Hess$ ha solo il contributo di $\A$ (curvatura $0$ dalla penalità,
perché $|w_i|$ è affine lì), mentre $\bm B$ ha in più il termine $\lam/|w^*_i|$, che è
\emph{grande} quando $|w^*_i|$ è piccolo. Quindi
\[
\frac{\bm v^{\!\top}\Hess\,\bm v}{\bm v^{\!\top}\bm B\,\bm v}
\approx\frac{\bm v^{\!\top}\A\,\bm v}{\bm v^{\!\top}\A\,\bm v+\lam\sum_i v_i^2/|w^*_i|}
\xrightarrow[\ |w^*_i|\to0\ ]{}0 .
\]

\begin{keybox}
\textbf{Meccanismo della lentezza.} IRLS modella $|w_i|$ con una parabola. Sulle componenti
attive $|w_i|$ è una retta (curvatura nulla), ma il modello quadratico $\bm B$ ci mette
curvatura $\lam/|w^*_i|$: \emph{tanta} dove la componente è piccola. Questo disallineamento
tra modello e funzione vera spinge $\mu_{\min}\to0$ e quindi $\spec\to1$. In una frase:
\emph{IRLS è lento proprio dove il LASSO è lineare}. Empiricamente $\spec\approx0.99$, e
quasi indipendente da $\eps$.
\end{keybox}

\noindent Caso limite opposto, utile per intuizione: se fosse $\bm B=\Hess$ avremmo
$\nabla\Mmap=\Id-\Id=\bm0$, rate $0$, convergenza in un passo (è ciò che fa il metodo di
Newton). IRLS è lento esattamente nella misura in cui il suo modello $\bm B$ si discosta dalla
vera Hessiana $\Hess$.

% ===================================================================
\section{Perché la figura sembra ``non lineare''}
\label{sec:figure}

Su un semilog del gap rispetto al LASSO esatto $f^\star$ contro $k$, la curva \emph{non} è una
singola retta. Si sovrappongono tre regimi:
\begin{enumerate}[nosep]
\item \textbf{Transiente iniziale veloce}: le prime iterazioni chiudono la maggior parte del
  gap (fase pre-asintotica, tratto ripido).
\item \textbf{Coda lineare lenta}: subentra il rate vero $\spec\approx0.99$, una retta quasi
  orizzontale; per vederla come retta servono \emph{molte} iterazioni.
\item \textbf{Pavimento di smussatura}: il gap rispetto a $f^\star$ non scende sotto
  $O(\eps)\approx10^{-8}$; con poche iterazioni la curva sbatte sul pavimento \emph{prima}
  che la coda lineare sia visibile.
\end{enumerate}
Sommati, danno una curva concava che a occhio sembra sublineare. La verifica corretta è
misurare il gap rispetto al minimo \emph{smussato} $f_\eps^\star$ (non $f^\star$), su molte
iterazioni: lì il rapporto per-iterazione $g_{k+1}/g_k$ \emph{si stabilizza} attorno a
$\approx0.995$ e oscilla senza crescere $\Rightarrow$ lineare, solo lento. Una sublineare
avrebbe $g_{k+1}/g_k$ in crescita monotona verso $1$.

% ===================================================================
\section{Avvertenze critiche (limiti della dimostrazione)}
\label{sec:caveats}

Da tenere a mente, anche all'orale:
\begin{enumerate}
\item \textbf{Rate locale, non globale.} Ostrowski garantisce attrazione e rate lineare in un
  \emph{intorno} di $\wstar$. La convergenza \emph{globale} della successione si argomenta a
  parte (monotonia MM $+$ forte convessità $+$ coercività $\Rightarrow$ ogni punto di
  accumulazione è il minimo unico). Sono due affermazioni distinte.
\item \textbf{Ipotesi ``nessuna componente sul bordo'' $|w^*_i|=\eps$.} Necessaria per la
  $C^1$-regolarità di $\Mmap$; è generica ma assunta, non dimostrata sul nostro problema.
\item \textbf{$\spec\approx0.99$ misurato su una sola istanza.} Per robustezza andrebbe
  confermato su piu' dataset/seed; la (quasi) indipendenza da $\eps$ è verificata, la
  variabilità tra istanze no. Ridurre il claim oppure rieseguire.
\item \textbf{Citazione da non usare.} Il rate lineare \emph{non} segue dal Teorema~5.3 di
  Daubechies et al.\ (quello assume la \emph{null space property} per sparse recovery
  \emph{sottodeterminato}): il nostro regime è \emph{sovradeterminato} e fortemente convesso.
  La giustificazione corretta è quella di questo documento: forte convessità del problema
  smussato $+$ teoria locale standard MM/Ostrowski.
\item \textbf{Gap relativi.} Per i confronti numerici usare gap \emph{relativi}
  ($|f-f^\star|/|f^\star|$): un gap assoluto di $10^{-6}$ è minuscolo se $f^\star\sim10^4$ e
  enorme se $f^\star\sim10^{-3}$.
\end{enumerate}

% ===================================================================
\section*{Riepilogo in cinque frasi (per l'orale)}
\addcontentsline{toc}{section}{Riepilogo in cinque frasi}
\begin{enumerate}[leftmargin=*]
\item IRLS è un'iterazione a punto fisso $\Mmap(\w)=\Q(\w)^{-1}\bvec$, MM su un problema
  smussato $f_\eps$ (Huber sul termine $L_1$), fortemente convesso perché $\X$ ha rango pieno.
\item Differenziando la mappa nel punto fisso si ottiene
  $\nabla\Mmap(\wstar)=\Id-\Q(\wstar)^{-1}\nabla^2 f_\eps(\wstar)$, la matrice di contrazione MM.
\item $\Q(\wstar)$ e $\nabla^2 f_\eps(\wstar)$ sono SPD e la loro differenza è
  $\lam\,\diag(1/|w^*_i|)$ sulle attive, $0$ sulle pinned, quindi $\succeq0$.
\item Ne segue che gli autovalori di $\Q^{-1}\nabla^2 f_\eps$ stanno in $(0,1]$, quelli di
  $\nabla\Mmap$ in $[0,1)$, dunque $\spec(\nabla\Mmap(\wstar))=1-\mu_{\min}<1$: per Ostrowski,
  convergenza lineare locale.
\item Il rate è vicino a $1$ perché sulle componenti attive $|w_i|$ ha curvatura nulla mentre
  il modello quadratico no: $\mu_{\min}\to0$ quando $|w^*_i|\to0$. Quindi IRLS converge
  linearmente ma lentamente; la figura sembra sublineare per transiente $+$ pavimento di
  smussatura, non perché lo sia.
\end{enumerate}

\end{document}
